\section{Residual regularity in radial singularity types}
\label{sec:radial-residual-regularity}

Subtracting a prescribed singularity separates regularity of the residual from the pole of the solution. On projective space, central radial mass gives an exact inverse for the equation. A relative density error then controls the residual through two derivatives, even when the density is merely integrable at the pole. The inverse also permits smooth crossing of density zeros away from it. The pole is the one used in \cref{lem:projective-sublogarithmic-density}; here its ordinary type is fixed while the density varies.

\subsection{The prescribed ordinary type}
Let $k\geq2$, let $\kappa$ be the Fubini--Study form on $\mathbb P^k$ with $\int\kappa^k=1$, and put $o=[1:0:\cdots:0]$. On the chart $Z_0\ne0$ write $z_i=Z_i/Z_0$ and set
\[
 r=\sum_{j=1}^k|z_j|^2,\qquad P=\frac r{1+r},\qquad
 T=1-\log P,
 \qquad u_\alpha=1-T^\alpha\quad(0<\alpha<1).
\]
The ratio $P=\sum_{i=1}^k|Z_i|^2/\sum_{i=0}^k|Z_i|^2$ is global. Give $u_\alpha$ its canonical value $-\infty$ at $o$ and value zero on the hyperplane $Z_0=0$. Primes denote differentiation in $r$. Throughout this section $b_\alpha=k(1-\alpha)>0$ and, near $r=0$, $L=1-\log r$.

\Needspace{12\baselineskip}
\begin{lemma}\label{lem:radial-nonlogarithmic-type}\index{Singularity type!radial nonlogarithmic pole}
The canonical $u_\alpha$ is $\kappa$-plurisubharmonic and smooth off $o$. Its non-pluripolar measure has a positive density $q_\alpha$ of mass one, in $L^1$ and in no $L^p$, $p>1$. Distinct $\alpha$ give distinct ordinary types, although all their multiplier ideals are trivial.
\end{lemma}
\begin{proof}
Plurisubharmonicity, smoothness off the pole, unit mass and the $L^p$ assertions are those of \cref{lem:projective-sublogarithmic-density}. To fix the profile used in the inverse below, put
\begin{equation}\label{eq:radial-residual-baseline}
 \Phi_\alpha=\log(1+r)+u_\alpha,\qquad
 Q_\alpha=r\Phi_\alpha'=P+\alpha(1-P)T^{\alpha-1}.
\end{equation}
The tangential and radial eigenvalues are $Q_\alpha/r$ and $Q_\alpha'$, both positive. Radial determinant integration, with the constant fixed by $\Phi=\log(1+r)$, gives
\begin{equation}\label{eq:radial-residual-flux}
 \int_{a<|z|^2<b}(\ddc\Phi)^k=Q(b)^k-Q(a)^k,
 \qquad Q=r\Phi'.
\end{equation}
Since $Q_\alpha(0)=0$ and $Q_\alpha(\infty)=1$, the density is
\begin{equation}\label{eq:radial-residual-baseline-density}
 q_\alpha=\frac{(Q_\alpha^k)'}{(P^k)'}>0.
\end{equation}
Smoothness on the hyperplane and the exclusion of the canonical
pole in the product leave no further measure. Near $r=0$, direct differentiation gives
\begin{align}\label{eq:radial-residual-asymptotics}
 Q_\alpha&\sim\alpha L^{\alpha-1},&
 Q_\alpha'&\sim\frac{\alpha(1-\alpha)}rL^{\alpha-2},\\
 (Q_\alpha^k)'&\sim
 \frac{k\alpha^k(1-\alpha)}rL^{-1-k(1-\alpha)}.\nonumber
\end{align}
In particular $q_\alpha\sim\alpha^k(1-\alpha)r^{-k}L^{-1-b_\alpha}$; the radial volume has order $r^{k-1}\,\mathrm{d}r$, recovering these integrability assertions.

The types for different $\alpha$ are ordinarily distinct, since
their difference is unbounded at $o$. Yet every $\mathcal I(tu_\alpha)$
is trivial for finite $t>0$: integrability of the constant function
reduces near $o$ to
$\int^\infty \mathrm{e}^{-kL+tL^\alpha}\,\mathrm{d}L<\infty$.
Ordinary type information here is strictly finer than this ideal
information.
\end{proof}

\Needspace{14\baselineskip}
\subsection{Density criterion and exact inverse}
Here ``radial'' means depending only on $r$ in the displayed chart.
The smoothness assumption along the hyperplane includes a smooth profile in $t=1/r$ near $t=0$.

\begin{theorem}\label{thm:radial-residual-regularity}\index{Regularity!radial singular residual}
Let $q\ge0$ be radial, smooth on $\mathbb P^k\setminus\{o\}$, including
the hyperplane at infinity, and satisfy $\int q\kappa^k=1$.
Suppose that for some $\delta,A,r_0>0$,
\begin{equation}\label{eq:radial-residual-density-error}
 \left|\frac{q(r)}{q_\alpha(r)}-1\right|\le Ar^\delta
 \quad(0<r<r_0).
\end{equation}
Then there is a radial $\kappa$-psh potential $U$, normalized
to have value zero on the hyperplane, with
\[
 \mathrm{MA}_\kappa(U)=q\kappa^k,
 \qquad U\sim u_\alpha.
\]
Its residual $R=U-u_\alpha$ off $o$ extends continuously
across $o$, with value $R_0$, and for small $r$ satisfies
\begin{equation}\label{eq:radial-residual-estimates}
 \begin{aligned}
 |R-R_0|&\le Cr^\delta L^{\alpha-2},\\
 |R'|&\le Cr^{\delta-1}L^{\alpha-2},\qquad
 |R''|\le Cr^{\delta-2}L^{\alpha-2}.
 \end{aligned}
\end{equation}
It is locally $C^{0,\min\{1,2\delta\}}$.
For $\delta\ge\frac12$ it extends $C^1$ with zero gradient at $o$;
for $\delta\ge1$ it extends $C^2$ with zero Hessian at $o$.
No derivative hypothesis on the error in \eqref{eq:radial-residual-density-error} is needed
for these conclusions.
The potential and its residual are smooth on $\mathbb P^k\setminus\{o\}$,
including every density zero there. Condition
\eqref{eq:radial-residual-density-error} makes the central mass of every ball
around $o$ positive; this positivity, not the smoothness of the density alone,
is what the smoothness across density zeros uses.
\end{theorem}

\begin{proof}
Let
\begin{equation}\label{eq:radial-residual-inverse}
 \begin{aligned}
 M(r)&=\int_{|z|^2<r}q\kappa^k,& Q(r)&=M(r)^{1/k},\\
 U'(r)&=\frac{Q(r)-P(r)}r,& U(\infty)&=0.
 \end{aligned}
\end{equation}
Condition \eqref{eq:radial-residual-density-error} makes $q$ positive near $o$, so $M(r)>0$ for every
$r>0$. One has $M'=q(P^k)'\ge0$, $Q'\ge0$ and $M(\infty)=1$.
Thus the positive root defining $Q$ is smooth for every finite $r>0$,
even where $q$ vanishes. This proves smoothness across all finite density
zeros without a vanishing-order assumption.
In $t=1/r$, the tail $1-M$ is the integral from $0$ to $t$ of the smooth density profile times $-\mathrm{d}(1+t)^{-k}/\mathrm{d}t$. Thus $M$ and $Q$ are smooth in $t$ near $t=0$, with $Q=1$ at $t=0$. The equation
$\mathrm{d}U/\mathrm{d}t=-(Q-P)/t$ is smooth there, since $Q-P$ vanishes at
$t=0$. This both defines the normalization and proves smoothness
along the full hyperplane.

From \eqref{eq:radial-residual-density-error} and
\eqref{eq:radial-residual-asymptotics},
\begin{equation}\label{eq:radial-residual-tail}
 |M-Q_\alpha^k|
 \le C\int_0^r s^{\delta-1}(1-\log s)^{-1-b_\alpha}\,\mathrm{d}s
 \le Cr^\delta L^{-1-b_\alpha}.
\end{equation}
The last inequality follows by writing $s=r\mathrm{e}^{-y}$ and bounding
$\int_0^\infty \mathrm{e}^{-\delta y}(L+y)^{-1-b_\alpha}\,\mathrm{d}y$.
Dividing by $Q_\alpha^k\asymp L^{-b_\alpha}$ shows that $Q/Q_\alpha$
tends to one; the mean value theorem for the positive $k$th root
then gives
\begin{equation}\label{eq:radial-residual-flux-error}
 |Q-Q_\alpha|\le Cr^\delta L^{\alpha-2}.
\end{equation}
Moreover the exact identity
\[
 \frac{Q'}{Q_\alpha'}=
 \frac q{q_\alpha}\left(\frac{Q_\alpha}{Q}\right)^{k-1}
\]
gives $|Q'-Q_\alpha'|\le Cr^{\delta-1}L^{\alpha-2}$.
Since $R'=(Q-Q_\alpha)/r$, integration to the origin gives a finite
$R_0$ and the first two inequalities in \eqref{eq:radial-residual-estimates}.
Differentiating this formula gives the last inequality.

The local potential $\Phi=\log(1+r)+U$ has semipositive
Hessian away from $o$ because $Q/r>0$ and $Q'\ge0$. Also
$U=u_\alpha+R_0+o(1)$ near $o$, so $\Phi$ tends to $-\infty$.
As a function of $\log r$, the profile of $\Phi$ has the nonnegative increasing slope $Q$, so the maxima $\max\{\Phi,\Phi(\varepsilon)\}$ are plurisubharmonic. They decrease to the locally integrable function $\Phi$ as $\varepsilon\downarrow0$, and the local decreasing-limit property shows that $\Phi$, with the value $-\infty$ at $o$, is plurisubharmonic. Thus $U$ is globally
$\kappa$-psh. Formula~\eqref{eq:radial-residual-flux} or the local determinant
ratio gives its measure $q\kappa^k$ away from $o$; the defining
non-pluripolar truncations exclude $o$ and are locally identical
to this smooth product elsewhere. The measure has exactly total
mass one. The bounded global residual proves the ordinary type.

In real coordinates, with $\rho=|z|=\sqrt r$,
\[
 |\nabla R|\le C\rho^{2\delta-1}L^{\alpha-2},
 \qquad |\nabla^2R|\le C\rho^{2\delta-2}L^{\alpha-2}.
\]
For $2\delta<1$, integration along radii bounds differences
by $C|\rho_1-\rho_2|^{2\delta}$ and hence by
$C|z_1-z_2|^{2\delta}$. For $2\delta\ge1$ the gradient is bounded,
giving local Lipschitz regularity. At $\delta=\frac12$ its bound
tends to zero and $R-R_0=o(\rho)$, which give the claimed $C^1$
extension. At $\delta=1$, the Hessian bound tends to zero and
$\nabla R=o(\rho)$, giving the $C^2$ extension. Larger $\delta$
only improves these bounds. These arguments justify crossing the
pole for $R$; they do not assign a finite value to $U(o)$.
\end{proof}

Among radial solutions with this density that are smooth off $o$,
including the hyperplane, and have no central flux ($Q(r)\to0$ as $r\to0$),
the potential is unique up to constants: its flux satisfies
$Q^k=M$, and \eqref{eq:radial-residual-inverse} then fixes its derivative.
No nonsymmetric uniqueness is asserted.

Since $U\sim u_\alpha$, \cref{prop:relative-elementary} gives relative envelope-class membership; an explicit stretching family shows it directly. If
$a_* =\inf R$, the $\kappa$-psh functions
$v_t=u_\alpha+t a_*$ satisfy, for every $t\ge1$,
\[
 t^{-1}v_t+(1-t^{-1})u_\alpha=u_\alpha+a_*\le U.
\]
At $t=1$ the term with zero coefficient is omitted, as in \cref{def:relative-envelope-class}. Hence
$U\in \mathcal E_{\mathrm{env}}(\mathbb P^k,\kappa,u_\alpha)$; at the common
pole $o$ both sides of the inequality equal $-\infty$.

\subsection{Smooth residuals from flat errors}
\begin{proposition}\label{prop:radial-flat-residual}
Under the hypotheses of \cref{thm:radial-residual-regularity}, suppose
$g=q/q_\alpha-1$ satisfies
\[
 (r\partial_r)^jg=O(r^N)
 \quad(r\to0)\quad\text{for every }N,j\ge0.
\]
Then $R-R_0$ extends smoothly across $o$, with every derivative
zero there.
\end{proposition}
\begin{proof}
Put $x=-\log r$, $a(x)=Q_\alpha(\mathrm{e}^{-x})^k$ and
$E(x)=M(\mathrm{e}^{-x})-a(x)$. The exact central integral is
\[
 E(x)=\int_x^\infty g(\mathrm{e}^{-y})[-a'(y)]\,\mathrm{d}y,
 \qquad E'(x)=g(\mathrm{e}^{-x})a'(x).
\]
Every derivative of $a$ has at most polynomial growth in $1+x$,
and $a\asymp(1+x)^{-b_\alpha}$. The integral bounds $E$ by an
exponential of arbitrarily large rate. Repeated differentiation of
$E'$ expresses each derivative as a finite Leibniz sum of flat
derivatives of $g$ and polynomially bounded derivatives of $a$;
all derivatives of $E$ are therefore exponentially flat.
Every derivative of $1/a$ has at most polynomial growth, by
\eqref{eq:radial-residual-baseline}. Thus $E/a$ has the same flatness. Applying
smooth composition near one to
\[
 Q-Q_\alpha=Q_\alpha\bigl[(1+E/a)^{1/k}-1\bigr]
\]
proves flatness of every $x$ derivative of $Q-Q_\alpha$.
Since $\mathrm{d}R/\mathrm{d}x=-(Q-Q_\alpha)$, integration from $x$ to infinity
proves the same for $R-R_0$, including its zeroth derivative.
An ordinary $j$th derivative in $r$ is $r^{-j}$ times a polynomial
in $\partial_x$; arbitrarily strong exponential decay absorbs that
factor. All ordinary derivatives extend continuously as zero at
$r=0$, and repeated one-dimensional integration verifies smoothness
there. Composition with $r=|z|^2$ gives the conclusion.
\end{proof}

\Needspace{19\baselineskip}
\subsection{Examples with sharp derivative thresholds}
\begin{proposition}\label{prop:radial-residual-sharpness}
The $C^1$ threshold $\delta=1/2$ and the $C^2$ threshold $\delta=1$ of \cref{thm:radial-residual-regularity} are sharp. When $2\delta<1$, the stated H\"older exponent is also sharp. There are radial probability densities $q\ne q_\alpha$ to which \cref{prop:radial-flat-residual} applies, so that $R-R_0$ is smooth and flat at $o$.
\end{proposition}
\begin{proof}
Choose smooth nonnegative radial cutoffs $\chi_0,\chi_1$ with
disjoint supports: $\chi_0=1$ near $r=0$ and vanishes outside a
finite small ball; $\chi_1\not\equiv0$ is supported in a finite annulus. Put
\begin{equation}\label{eq:radial-residual-sharp-examples}
 B=\frac{\int r^\delta\chi_0q_\alpha\kappa^k}
          {\int\chi_1q_\alpha\kappa^k},\qquad
 q_c=q_\alpha\bigl[1+c(r^\delta\chi_0-B\chi_1)\bigr].
\end{equation}
For sufficiently small $c>0$, $q_c$ is positive, smooth off $o$
and along the hyperplane, and has total mass exactly one. It is
comparable to $q_\alpha$ near $o$, hence is in no $L^p$, $p>1$.
\Cref{thm:radial-residual-regularity} gives a solution on $\mathbb P^k$ in
the prescribed ordinary type. Near $o$ the difference $M_c-Q_\alpha^k$ of the central masses is
exactly $c\int_0^r s^\delta(Q_\alpha^k)'\,\mathrm{d}s$.
The substitution $s=r\mathrm{e}^{-y}$ and dominated convergence give
\[
 \int_0^r s^{\delta-1}(1-\log s)^{-1-b_\alpha}\,\mathrm{d}s
 \sim\delta^{-1}r^\delta L^{-1-b_\alpha}.
\]
Using \eqref{eq:radial-residual-asymptotics} and the $k$th-root expansion yields
\begin{align*}
 R_c'(r)&\sim\frac{c\alpha(1-\alpha)}\delta
                  r^{\delta-1}L^{\alpha-2},\\
 R_c(r)-R_c(0)&\sim\frac{c\alpha(1-\alpha)}{\delta^2}
                  r^\delta L^{\alpha-2}.
\end{align*}
For $\delta<\frac12$, the radial gradient
$2\sqrt r\,R_c'$ diverges, ruling out $C^1$ extension. For
$\delta<1$, a tangential real Hessian eigenvalue $2R_c'$ diverges,
ruling out $C^2$ extension. The annular compensation does not change the
central mass near $o$. For $2\delta<1$ the displayed asymptotic also rules
out any larger H\"older exponent. The critical derivative thresholds
in \cref{thm:radial-residual-regularity} are therefore sharp for this family.

Replacing $r^\delta\chi_0$ in \eqref{eq:radial-residual-sharp-examples} by
$\mathrm{e}^{-1/r}\chi_0$ and adjusting $B$ by the same integral formula
gives a nonzero mass-preserving perturbation satisfying every
flat-error hypothesis. \Cref{prop:radial-flat-residual} then shows that
$R-R_0$ is smooth and flat at the pole.
\end{proof}

\subsection{Smooth residuals with an open density-zero annulus}
\begin{proposition}\label{prop:radial-zero-annulus}
For every $k\ge2$ and $0<\alpha<1$, there is a nonnegative radial
probability density $q_0$, smooth off $o$ and along the hyperplane,
which equals $q_\alpha$ near $o$ and vanishes on a nonempty open
annulus. Its solution has ordinary type $u_\alpha$ and a
globally smooth residual, constant near $o$.
\end{proposition}
\begin{proof}
Choose a smooth radial $h$ with $0\le h\le1$, equal to zero on
a nonempty finite annulus and equal to one near $o$ and outside
a larger finite annulus. Choose a nonzero nonnegative smooth radial
cutoff $\chi$ supported in a disjoint finite annulus where $h=1$.
Set
\[
 C=\frac{\int(1-h)q_\alpha\kappa^k}{\int\chi q_\alpha\kappa^k}>0,
 \qquad q_0=q_\alpha(h+C\chi).
\]
This is nonnegative and has integral one. It is smooth at every
point off $o$, including the hyperplane, and equals $q_\alpha$
near $o$. It vanishes on the chosen open annulus and is still in
no $L^p$, $p>1$. \Cref{thm:radial-residual-regularity} applies, with
zero relative density error on a neighborhood of $o$. There the
exact central fluxes agree, so $Q=Q_\alpha$ and $R'=0$.
Hence $R$ is constant near $o$ and extends smoothly across it.
It is smooth everywhere else by the positive-root inverse, including
the entire zero annulus. The solution itself keeps its pole at $o$;
the smoothness assertion concerns the residual.
\end{proof}

\subsection{Products with a zero-lower-mass factor}
Let $(X_0,h_0)$ be any connected compact Hermitian $m$-fold with
$v_-(h_0)=0$, and write $V_0=\int_{X_0}h_0^m>0$.
On $X=X_0\times\mathbb P^k$ put $N=m+k$ and
$H=\operatorname{pr}_0^*h_0+\operatorname{pr}_Y^*\kappa$, where $\operatorname{pr}_0$ and $\operatorname{pr}_Y$ are the projections onto $X_0$ and $Y=\mathbb P^k$. Let $q$ and $U$ be as in \cref{thm:radial-residual-regularity}; in the product formulas below, $U$ and $R$ denote their pullbacks.
\begin{corollary}\label{cor:radial-residual-nonclosed-product}
The pullback of $U$ has the ordinary type of the pullback of $u_\alpha$, belongs to its relative envelope class $\mathcal E_{\mathrm{env}}(X,H,u_\alpha)$ and has finite Monge--Amp\`ere mass $\binom Nm V_0$. Its residual has the regularity of \cref{thm:radial-residual-regularity} across $X_0\times\{o\}$. The background $H$ is nonclosed and $v_-(H)=0$.
\end{corollary}
\begin{proof}
The pullback of $U$ obeys
\[
 \mathrm{MA}_H(U)=\binom Nm\operatorname{pr}_0^*(h_0^m)
                 \wedge\operatorname{pr}_Y^*(q\kappa^k),
 \qquad \int_X\mathrm{MA}_H(U)=\binom Nm V_0.
\]
Indeed the bounded truncation $\max\{U,-j\}$ is a pullback.
The bounded corrected-product expansion of \cref{prop:bounded-products,lem:bounded-products-submersion} leaves only the displayed
binomial term by dimensions; multiplication by $\mathbf1_{\{U>-j\}}$
and increasing $j$ give the exact formula and finite mass.
The residual is the pullback of $R$, so its regularity
extends across $X_0\times\{o\}$. It is ordinarily equivalent to the
pullback of $u_\alpha$ and belongs to the corresponding relative
stretching class by the same explicit $v_t$.

The last two assertions are also part of \cref{thm:zero-lower-mass-product-equation}. Directly,
$v_-(H)=0$ because bounded $h_0$-psh potentials with masses tending
to zero have pullbacks of mass $\binom Nm$ times those masses.
The background is nonclosed, since a closed positive
$h_0$ would have all bounded masses equal to $V_0>0$ by Stokes.
The conclusion $v_-(H)=0$ uses the hypothesis $v_-(h_0)=0$ on the factor, and this calculation does not determine the threshold of an individual density.
\end{proof}
